Suposem un llamp núvol-terra negatiu (una transferència de càrrega negativa del núvol cap a terra) que passa per un parallamps (barra metàl·lica vertical). En aquest moment es crea un camp magnètic al voltant del parallamps que podem modelar al creat per un fil de corrent infinit.

\begin{apartat}{1,25}
Calculeu el camp magnètic a una distància de 10~cm del parallamps quan hi passa un corrent de 100~kA. Feu un dibuix esquemàtic del parallamps indicant el sentit del moviment dels electrons, la intensitat de corrent i tres línies de camp magnètic. Justifiqueu el sentit de les línies de camp.
\end{apartat}

\begin{apartat}{1,25}
Representeu gràficament el mòdul d'aquest camp magnètic en funció de la distància $r$ al parallamps en un interval de $10\ \mathrm{cm} \le r \le 50\ \mathrm{cm}$. Suposem un electró que en el moment de la descàrrega es troba a 10~cm del parallamps i amb una velocitat de $10^{3}\ \mathrm{m/s}$ paral·lela al parallamps i cap a terra. Calculeu el mòdul de la força que crea el camp magnètic sobre l'electró i justifiqueu quins en són la direcció i el sentit.
\end{apartat}

\dades{%
  El mòdul del camp magnètic creat per un fil infinit per on circula un corrent $I$ a una distància $r$ del fil és $B = \dfrac{\mu_0 I}{2\pi r}$.\\[4pt]
  $\mu_0 = 4\pi\times 10^{-7}\ \mathrm{T\cdot m\cdot A^{-1}}$\\
  $|e| = 1,602\times 10^{-19}\ \mathrm{C}$}

\figura[0.55]{fig1.png}
